% ================================================
% 文件: 03_通信接收机.tex
% 章节: 第8章 通信接收机
% 所属部分: 接收机技术
% 版本: v1.0
% ================================================
% 本章目录结构（树状）:
% \chapter{通信接收机}
% ├── \section{通信接收机概述}
% ├── \section{短波通信接收机}
% ├── \section{单边带接收机}
% ├── \section{通信接收机设计要点}
% └── \section{通信接收机测试与校准}
% ================================================

\chapter{通信接收机}
\label{chap:communication_receiver}

\section{通信接收机概述}
\label{sec:comm_receiver_intro}

通信接收机是用于专业通信的无线电接收机，
主要用于点对点通信。
与广播接收机相比，
通信接收机具有以下特点：

\begin{itemize}
    \item \textbf{宽频段覆盖}：通常覆盖从低频到微波的多个频段
    \item \textbf{多种调制方式}：支持AM、FM、SSB、CW等多种调制
    \item \textbf{高灵敏度}：能够接收微弱信号
    \item \textbf{良好的选择性}：能够区分邻近信道
    \item \textbf{稳定性要求高}：频率稳定度和温度稳定性要求严格
\end{itemize}

专业通信接收机与家用广播接收机在设计目标上有本质区别。
广播接收机以“听得清楚、
操作简单”为首要目标，
往往只覆盖本地广播频段；
而通信接收机面向固定台站、
船舶、
航空、
应急以及业余无线电等用途，
要求在一台设备内连续覆盖从长波、
中波、
短波一直到甚高频（VHF）与特高频（UHF）的宽频带，
并能适应不同时期、
不同业务的调制方式。
为此，
通信接收机普遍采用超外差（superheterodyne）架构，
并配合现代频率合成技术与数字信号处理（DSP），
在很宽的输入电平范围内保持稳定的接收性能。

所谓超外差，
是利用本地振荡信号与输入射频信号在混频器中相乘，
产生一个频率固定的中频信号，
后续的中频放大、
滤波与解调均围绕这个固定中频展开。
其核心关系式为中频频率等于射频频率与本地振荡频率之差的绝对值：

\[
f_{\mathrm{IF}}=\left| f_{\mathrm{LO}}-f_{\mathrm{RF}} 
\right|
\]

其中 $f_{\mathrm{RF}}$ 为接收的射频频率，$f_{\mathrm{LO}}$ 为本地振荡频率，$f_{\mathrm{IF}}$ 为变换后的中频。由于中频固定，接收机可以在这一频率上获得很高的增益与陡峭的选择性，而不必像直接放大式接收机那样在每个波段都重新设计调谐放大器。这一思想（详见第\ref{chap:receiver_fundamentals}章“接收机基础”）是现代通信接收机的共同基础。

现代通信接收机还强调“可维护性”与“可测量性”：
其内部增益、
带宽、
解调模式、
AGC（自动增益控制）特性往往可由操作员或测试程序设定，
从而满足不同链路、
不同业务的工程需要。
本章后续各节将围绕短波接收、
单边带解调、
工程设计要点与测试校准展开讨论。

接收机的灵敏度须放在“链路预算”中理解：发射功率、天线增益、传播损耗与接收机灵敏度共同决定可通距离。热力学给出噪声底 $kTB$（$k$ 为玻尔兹曼常数、$T$ 为温度、$B$ 为带宽），任何接收机都无法低于此底；噪声系数只是描述接收机“比理想多引入多少噪声”。在卫星通信等场景中，常用品质因数 $G/T$（天线增益与系统噪声温度之比）而非单纯噪声系数衡量整体接收能力，因为大天线增益同样关键。

现代通信接收机越来越多地融入软件定义无线电（SDR）生态：
以通用射频前端加算力平台运行 GNURadio、
MATLAB 等框架，
研究者与爱好者可快速验证新波形与新算法，
也推动专业设备向开放架构演进。
但无论架构如何变化，
灵敏度、
选择性、
动态范围与杂散这四项底层指标仍是评价任何通信接收机的根本，
软件只是把它们在不同程度上从硬件搬到了数字域实现。

\section{短波通信接收机}
\label{sec:shortwave_receiver}

短波通信接收机工作在3-30MHz频段，
利用电离层反射进行远距离通信。

特点：
\begin{itemize}
    \item \textbf{天波传播}：信号经电离层反射，可实现数千公里通信
    \item \textbf{信道拥挤}：频段资源有限，信道密集
    \item \textbf{信号衰落}：受电离层变化影响，信号强度波动大
    \item \textbf{多径传播}：同一信号经多条路径到达，产生衰落
\end{itemize}

短波接收机通常采用超外差架构，
具有高灵敏度和良好的选择性。

短波波段（3\sim30MHz）的无线电波主要依靠电离层中的 D、
E、
F 层反射实现天波传播，
单次反射即可跨越数千公里，
经多次“跳越”（hop）甚至可环绕半球。
由于电离层密度随昼夜、
季节、
太阳活动（黑子数、
耀斑）剧烈变化，
可用的最高可用频率（MUF）与最低可用频率（LUF）也在不断变化，
因此短波接收机必须能够在整个波段内连续调谐，
并具备足够的灵敏度余量以应付深衰落。

在短波信道中，
信号衰落表现为幅度与相位的随机起伏。
多径传播是衰落的主要来源：
同一发射信号经不同反射路径（或多跳）到达接收端，
各路径时延不同，
在接收点叠加后有时同相增强、
有时反相抵消，
形成频率选择性衰落；
当两条路径时延差接近一个信号周期时，
还会产生选择性衰落，
使某些频率分量被严重衰减。
为对抗衰落，
短波接收机通常设有噪声抑制（noise blanker）、
可调中频带宽（如 500Hz、
2.4kHz、
6kHz 可选）以及较快的 AGC 响应。

短波频段业务密集，
除业余无线电、
海事、
航空、
广播外，
还可能存在各种窄带与宽带干扰。
因此短波通信接收机几乎都采用二次变频（double conversion）方案：
第一中频取得较高的镜频抑制与较好的前端隔离，
第二中频（通常为 455kHz 或 450kHz）再完成陡峭的邻道选择。
其本振与混频链路由高稳定度的频率合成器驱动，
频率步进可做到 10Hz 甚至 1Hz，
以适应 CW、
SSB 等窄带模式的精细调谐。

现代短波接收机在模拟中频之后普遍引入 16\sim24 位高速 ADC 与 DSP，
完成数字滤波、
解调、
噪声消除与频谱显示，
显著改善了邻道选择性与可操作性。
需要指出的是，
短波接收机的性能上限仍受天线与传播条件制约，
再好的接收机也无法恢复被噪声完全淹没或受强邻台阻塞的信号。

\section{单边带接收机}
\label{sec:ssb_receiver}

单边带（SSB）是一种高效的调制方式，
只传输调制信号的一个边带。

SSB接收机的特点：
\begin{itemize}
    \item \textbf{带宽窄}：仅为AM信号的一半，节省频谱资源
    \item \textbf{功率效率高}：发射功率集中在有用信号上
    \item \textbf{需要载波恢复}：接收端需重建载波才能正确解调
    \item \textbf{对频率稳定度要求高}：频率偏差会导致信号失真
\end{itemize}

SSB广泛应用于短波通信、
业余无线电等领域。

单边带信号在发射端抑制了载波、
并滤除一个边带（上边带 USB 或下边带 LSB），
因此接收端收到的只是一个“搬移了频率的基带”，
本身不包含可直解的载波。
SSB 解调的核心是一个乘积检波器（product detector，
也称同步检波器）：
用接收机内部一个叫本机载波振荡器（BFO，
beat frequency oscillator）的信号与输入的 SSB 信号相乘，
将其边带重新搬回音频。
解调关系可表示为

\[
s_{\mathrm{out}}(t)=s_{\mathrm{SSB}}(t)\cdot 2\cos(2\pi f_{\mathrm{BFO}} t)
\]

当 BFO 频率 $f_{\mathrm{BFO}}$ 与发射端被抑制载波的频率精确一致时，输出即为原始音频；一旦存在频差 $\Delta f$，输出音频整体平移 $\Delta f$，使语音变调、难以辨认。工程上通常要求接收机与发射端的频差控制在数十赫兹以内（如 $\pm50$Hz），这正是 SSB 对频率稳定度要求高、且“需要载波恢复”的根本原因。

由于 SSB 没有载波可供 AGC 提取包络，
其自动增益控制需从信号本身的幅度起伏中提取，
故 SSB 的 AGC 时间常数通常取得比 AM 慢，
以避免随语音音节起伏产生“ pumping”失真。
同时，
SSB 接收机必须准确选择边带：
对 USB 信号若误用 LSB 边带滤波器，
解调出的将是一个颠倒、
不可懂的频谱，
因此操作员需正确设置边带与中心频率。

为获得良好的话音可懂度，
SSB 接收机的有效中频带宽一般取 2.1\sim2.8kHz（典型 2.4kHz），
并配合陡峭的晶体或机械滤波器以抑制邻道。
在弱信号条件下，
有时进一步压缩到 1.8kHz 甚至 500Hz 的 CW 带宽以提高信噪比。
SSB 接收机也常用于数字模式（如 RTTY、
PSK31、
FT8）的前端，
解调后的音频再交由计算机进行数字解码。

单边带接收机对解调电路的要求集中在“载波恢复”与“带宽整形”两点。
乘积检波器由乘法器或环形解调器构成，
BFO 的相位与频率需与发射抑制载波锁定；
为容许收发的微小频差，
现代设备常用 AFC 或数字载波跟踪环把 BFO 缓慢锁到信号谱中心。
AGC 在 SSB 中多采用“慢速”或“自适应”模式，
避免语音音节起伏被当作幅度信息放大而产生泵浦失真；
部分接收机还提供“快/慢/关”三挡供操作员选择。

除带宽外，
SSB 接收机常配备通带调谐（passband tuning）与可变中心频率，
使操作员能把干扰推出通带或补偿对方频偏；
独立的音频陷波（用于抑制特定单音干扰）和音频峰值滤波（用于 CW/窄带）也常见于高性能机型。
当 SSB 通路用于数字模式时，
要求线性相位与平坦幅频，
以免数据码间串扰；
此时常把中频带宽放宽到 3kHz 以上并关闭不必要的处理。
可见 SSB 接收机的“解调”并非单点电路，
而是一组围绕边带、
载波与带宽的协同处理。

\section{通信接收机设计要点}
\label{sec:comm_receiver_design}

设计通信接收机时需考虑：

\begin{itemize}
    \item \textbf{频率合成器}：提供稳定、准确的本地振荡频率
    \item \textbf{前端设计}：低噪声放大器、预选滤波器
    \item \textbf{中频处理}：高Q值滤波器、自动增益控制（AGC）
    \item \textbf{解调方式}：支持多种调制方式的解调电路
    \item \textbf{干扰抑制}：镜像抑制、邻道抑制
\end{itemize}

除上述基本环节外，
通信接收机的工程设计还需在“灵敏度、
选择性、
动态范围”三者之间反复权衡，
以下分述关键模块与指标。

\medskip\noindent\textbf{频率合成与前端的取舍}

频率合成器决定本振的准确度、
分辨力与相位噪声。
早期接收机采用 LC 振荡或晶体切换，
现代则普遍采用锁相环（PLL）甚至直接数字合成（DDS）方案，
可在全频段内获得 1Hz 级的分辨力与较低的相位噪声。
前端由预选滤波器（preselector）与低噪声放大器（LNA）组成：
预选器在混频之前滤除带外强信号，
降低镜频与互调风险；
LNA 则置于最前端以压低整机噪声系数，
但其增益过高又会压缩后端动态范围，
故 LNA 的增益与类型需与混频器、
后续增益统筹设计。

\medskip\noindent\textbf{一次变频与二次变频、高中频与低中频}

通信接收机可按变频次数分为一次变频与二次变频两类。一次变频架构简单、成本低，但第一中频若取得较低，则镜频距信号仅 $2f_{\mathrm{IF}}$，对前端预选要求苛刻；若第一中频取得较高（高中频方案），则镜频远离信号，镜频抑制容易实现，但高中频滤波器的矩形系数较难做陡。二次变频先用一个较高（或较低）的第一中频解决镜频/带宽问题，再用第二中频完成高选择性的邻道滤波，是高端通信接收机的主流。镜频频率（高本振注入，即 $f_{\mathrm{LO}}>f_{\mathrm{RF}}$ 时）为

\[
f_{\mathrm{image}}=f_{\mathrm{RF}}+2f_{\mathrm{IF}}
\]

可见第一中频越高，
镜频离信号越远，
越容易用预选器将其滤除。

\medskip\noindent\textbf{晶体滤波器、动态范围、阻塞与互调}

中频滤波器决定选择性，
其中石英晶体滤波器具有极高的 Q 值与矩形度，
是实现 2.4kHz、
500Hz 等窄带选择性的关键器件；
其缺点是中心频率固定，
通常只在第二中频使用。
接收机的动态范围描述其同时处理“微弱有用信号”与“强干扰”的能力，
常用“阻塞动态范围”（从最小可检测信号 MDS 到产生规定增益压缩或阻塞的强信号电平）或“无杂散动态范围（SFDR）”来表征。
阻塞（blocking）指强邻道信号使接收机增益下降、
灵敏度恶化的现象；
互调（intermodulation）则是两个或多个强信号经前端非线性混合后，
在有用频率上产生虚假响应，
最典型的是三阶互调。

整机噪声系数可由各级级联求出（Friis 公式）。
以两级为例：

\[
F=F_1+\frac{F_2-1}{G_1}
\]

其中 $F_1$ 为第一级（通常即 LNA）的噪声系数，$G_1$ 为其功率增益，$F_2$ 为第二级噪声系数。该式说明：只要第一级增益足够高、噪声系数足够低，后续各级对整机噪声的贡献就被“压缩”，因此 LNA 的位置与性能至关重要。

\medskip\noindent\textbf{AGC 与静噪}

AGC 根据接收信号强弱自动调节中频与射频增益，
使输出电平在输入变化数十 dB 时基本稳定；
其时间常数需依模式设定（SSB 慢、
FM 快）。
静噪（squelch）则用于在信号低于门限时关闭音频输出，
避免无信号时的噪声干扰，
常用于 FM 与语音信道。

为便于方案比较，
下表列出一次变频与二次变频接收机的主要特征。

\begin{table}[htbp]
\centering
\caption{通信接收机一次变频与二次变频方案比较}
\begin{tabular}{|p{4.5cm}|p{5cm}|p{5cm}|}
\hline
\textbf{比较项目} & \textbf{一次变频} & \textbf{二次变频} \\
\hline
结构复杂度 & 简单、成本低 & 较复杂、成本较高 \\
\hline
镜频抑制难度 & 中频低时难，需强预选 & 高中频下镜频远离，易实现 \\
\hline
邻道选择性 & 依赖单一中频滤波器 & 第二中频可用晶体滤波器，更陡 \\
\hline
适用场合 & 宽带扫描、低成本设备 & 高性能短波/专业接收机 \\
\hline
\end{tabular}
\end{table}

在工程设计层面，镜像抑制与中频选择是一对核心权衡。若采用一次变频且第一中频较低，镜频仅距信号 $2f_{\mathrm{IF}}$，前端预选器必须非常陡峭才能抑制；若把第一中频取高（如数十 MHz 量级），镜频远离、预选容易，但高中频上的晶体滤波器难以做到很陡，邻道选择性需靠第二中频完成。因此高端接收机多用“高中频一次变频 + 低中频二次变频”的组合：第一中频解决镜频与宽带，第二中频（常用 455kHz 或 450kHz）用晶体/机械滤波器解决选邻道。屏蔽与接地同样关键——本振泄漏、中频辐射都会造成自激或虚假响应，需通过分层屏蔽与滤波埠抑制。

杂散响应还与本振相位噪声、频率合成步进相关：本振近端相位噪声过大会抬升接收机噪声基底、恶化弱信号；合成步进过粗则不利于 CW/SSB 精调。设计者需在 PLL 环路带宽、参考频率与分频比之间折中，并对可能出现的互调产物（如 $m f_{\mathrm{LO}} \pm n f_{\mathrm{RF}}$ 落入中频）做频率规划规避。良好的接收机应在“灵敏度—选择性—动态范围—杂散”四维指标上取得平衡，而不是单点最优。

\section{通信接收机测试与校准}
\label{sec:comm_receiver_testing}

测试项目包括：

\begin{itemize}
    \item \textbf{灵敏度测试}：测量最小可检测信号
    \item \textbf{选择性测试}：测量邻道抑制能力
    \item \textbf{频率准确度}：测量频率误差
    \item \textbf{失真测试}：测量信号失真程度
    \item \textbf{AGC特性}：测量自动增益控制范围和响应速度
\end{itemize}

定期校准可确保接收机性能符合指标要求。

通信接收机的定量化测试以灵敏度和失真表征为主。最小可检测信号（MDS）通常定义为在给定中频带宽与噪声系数下，输出信纳比（SINAD）达到规定值（如 12dB）时所需的最小输入电平。热噪声底为 $-174$ dBm/Hz，故

\[
P_{\mathrm{MDS}}=-174+10\log_{10}(B)+NF+\left(\frac{S}{N+D}\right)_{\min}
\]

其中 $B$ 为中频噪声带宽（Hz），$NF$ 为整机噪声系数（dB），$(S/(N+D))_{\min}$ 为所需最小信纳比（dB）。SINAD 表示信号加噪声加失真与纯噪声加失真之比，常用分贝表示：

\[
\mathrm{SINAD}=20\log_{10}\frac{S+N+D}{N+D}
\]

当以 12dB SINAD 作为灵敏度判据时，上式最后一项即为 12dB。选择性测试则测量在偏离中心频率某一间隔（如 $\pm10$kHz、$\pm25$kHz）处，接收机对干扰的衰减量，反映中频滤波器的邻道抑制能力。

三阶互调截获点（IP3）是评价强信号下线性度的关键指标。向接收机注入两个幅度相等、频率相近的测试音 $f_1$、$f_2$，其互调产物出现在 $2f_1-f_2$ 与 $2f_2-f_1$。当输入电平提高 $\Delta$（dB）时，基波输出提高 $\Delta$，三阶产物提高 $3\Delta$；将二者外推至相等的交点即输入（或输出）三阶截获点。IP3 越高，接收机在强干扰环境中越不易产生互调虚假响应。

测试通常在射频屏蔽室内用标准信号发生器、
音频分析仪与功率计完成。
校准内容包括本振频率准确度（与频率标准比对）、
中频带宽与中心频率、
S 表/电平读数线性、
AGC 范围与响应时间。
下表汇总典型测试项目与常用判据。

现场测试还应关注“场强”与“相对电平”的换算。接收机面板的 S 表通常按 1 个 S 单位对应 6dB（或以 $-73$ dBm 为 S9 的惯例）标定，但不同机型刻度并不一致，精确测量仍需用标准信号发生器注入已知电平、绘制“输入电平—S 表读数”校准曲线。对于固定台站，还常以标准测试天线换算到场强（dB$\mu$V/m），以评估实际链路余量。

校准周期视使用环境而定，
移动与野外设备建议每次任务前做频率与灵敏度核查，
固定台站可每半年至一年做全面计量。
需要强调的是，
实验室给出的“灵敏度”是在无外部干扰、
标准温度与标准天线端口条件下的指标；
实际开路的灵敏度和可通率还受天线增益、
馈线损耗与传播条件强烈影响，
不能把实验室指标直接等同于现场性能。

计量溯源保证“测得准”：
灵敏度、
频率准确度、
IP3 等应由可溯源至国家标准的设备定期校准，
出具不确定度报告；
固定台站可建立自动测试系统，
按计划循环完成各指标自检并归档趋势，
提前发现增益漂移或前端劣化。
现场快速核查则侧重“能否工作”：
用已知电平信号验证 S 表、
用标准调制验证解调、
用频谱仪观察杂散，
足以在任务前排除明显故障。

日常使用中还应建立“频率—模式—带宽—天线”的参数档案，
使不同操作员获得一致设置；
对强干扰环境，
依靠预选、
衰减与杂散表管理接收机响应。
通信接收机虽以“通得了”为最终目标，
但只有把灵敏度、
选择性、
动态范围、
互调与校准都管理到位，
才能在复杂电磁环境中稳定、
可重复地达成这一目标。
规范测试与校准，
是专业通信接收机从“能响”走向“可信”的必经之路。

\begin{table}[htbp]
\centering
\caption{通信接收机典型测试项目与判据}
\begin{tabular}{|p{4cm}|p{5.5cm}|p{5cm}|}
\hline
\textbf{测试项目} & \textbf{测试方法} & \textbf{常用判据} \\
\hline
灵敏度（MDS） & 12dB SINAD 法 & 视带宽与 NF，通常优于 -120dBm \\
\hline
选择性 & 偏离中心测衰减 & 邻道衰减 > 60dB 为佳 \\
\hline
频率准确度 & 与标准频率比对 & 误差 < 1\times10^{-6}（温补） \\
\hline
三阶互调 IP3 & 双音法测截获点 & 越高越好，典型 +10\sim+30dBm \\
\hline
AGC 特性 & 改变输入测输出变化 & 输入变化 60dB，输出变化 < 6dB \\
\hline
\end{tabular}
\end{table}
